\providecommand{\Sp}{\operatorname{Sp}}
\let\HH\relax
\newcommand{\HH}{\mathbb{H}}

\chapter{James Arthur (2015)}
\index{Arthur, James}

\begin{quote}
\it ``James Arthur has transformed the theory of automorphic forms through his development of the trace formula and his classification of representations of classical groups. His work is a monumental achievement that has laid the foundation for much of modern number theory.'' --- Wolf Prize Citation, 2015
\end{quote}

\section{Early Life and Career}

James Arthur was born in 1944 in Hamilton, Ontario, Canada. He studied at the University of Toronto, receiving his Ph.D.\ in 1970 under the supervision of David Lowdenslager at Yale University. He held positions at Yale University and the University of Michigan before returning to the University of Toronto, where he has been a professor since 1978. He was awarded the Wolf Prize in Mathematics in 2015.

\section{Main Contributions}

James Arthur (born 1944) is a Canadian mathematician whose work on the trace formula and automorphic representations has fundamentally shaped modern number theory and representation theory. He was awarded the Wolf Prize in Mathematics in 2015.

The Wolf Prize citation reads: ``for his monumental work on the Selberg trace formula and on the representation theory of reductive groups, and for his fundamental contributions to the Langlands program.''

Arthur's core contributions include:
\begin{itemize}
    \item \textbf{Arthur--Selberg Trace Formula\index{trace formula, Arthur-Selberg}:} A far-reaching generalization of Selberg's trace formula that applies to any reductive group over a number field, relating automorphic representations to orbital integrals on the group.
    \item \textbf{Arthur's Conjectures:} A conjectural classification of the discrete automorphic spectrum of any reductive group in terms of global Langlands parameters, now a central part of the Langlands program.
    \item \textbf{Endoscopic Classification:} The proof (building on work of many others) of the endoscopic classification of representations of classical groups, establishing the Langlands functoriality principle for these groups.
    \item \textbf{The Arthur--Clozel Base Change:} A fundamental result establishing base change for $\GL(n)$ as predicted by Langlands functoriality.
    \item \textbf{Local Trace Formula:} The development of a local analogue of the trace formula, essential for the study of local Langlands correspondence.
\end{itemize}

\section{Origin of the Problem}

\subsection{Selberg's Trace Formula}

In 1956, Selberg proved a remarkable formula for a Riemannian surface $\Gamma \backslash \HH^2$:
\[
\sum_{\text{eigenvalues}} h(\lambda_j) = \sum_{\text{conjugacy classes } \{\gamma\}} \frac{\operatorname{vol}(\Gamma \backslash \HH^2)}{2\pi} \int_{-\infty}^{\infty} h(r) \tanh(\pi r) \, dr + \sum_{\gamma \neq 1} \frac{\ell(\gamma_0)}{2\pi \sinh(\ell(\gamma)/2)} \hat{h}(\ell(\gamma))
\]

This formula relates the spectral data (eigenvalues of the Laplacian) to geometric data (lengths of geodesics). Selberg's formula works for $\GL(2)$ (or $\SL(2,\R)$), but a generalization to higher rank groups was needed for the Langlands program.

\subsection{The Langlands Program}

Langlands proposed a grand unified vision connecting number theory, representation theory, and harmonic analysis. At its core is the principle of functoriality: for any morphism of Langlands dual groups $^L G \to ^L H$, there should be a transfer of automorphic representations from $G$ to $H$. Proving such transfers requires powerful analytic tools.

\subsection{The Need for a General Trace Formula}

To compare automorphic spectra of different groups (as required by functoriality), one needed a trace formula that:
\begin{enumerate}
    \item Works for any reductive group over a number field.
    \item Can be stabilized (grouped by stable conjugacy classes rather than rational conjugacy classes).
    \item Can be compared between different groups (the ``comparison of trace formulas'').
\end{enumerate}

\section{How It Was Solved and the Intuitions/Tools Needed}

\subsection{Arthur's Trace Formula}

Arthur developed the trace formula through a series of monumental papers from 1978 to 2000. The key steps include:
\begin{enumerate}
    \item \textbf{The Coarse Expansion:} Decomposing the kernel of the trace formula according to the cuspidal data and the unipotent radical.
    \item \textbf{The Fine Expansion:} Further decomposing each coarse term into contributions from individual automorphic representations.
    \item \textbf{The Invariant Trace Formula:} Removing the dependence on a choice of truncation parameter to obtain an invariant formula.
    \item \textbf{The Stable Trace Formula:} Refining the invariant formula to group terms by stable conjugacy classes.
\end{enumerate}

\subsection{Tools and Techniques}

\begin{itemize}
    \item \textbf{Geometric Analysis:} The spectral theory of the Laplacian and the analysis of Eisenstein series, building on the work of Langlands.
    \item \textbf{Harmonic Analysis on Reductive Groups:} The theory of orbital integrals, Harish-Chandra's Plancherel theorem, and the theory of characters.
    \item \textbf{Stable Conjugacy:} The notion that conjugacy over an algebraic closure, rather than over the base field, gives the correct geometric terms for comparison.
    \item \textbf{Endoscopy:} The theory of endoscopic groups, which are auxiliary groups whose trace formulas can be compared to the original group.
\end{itemize}

\section{Mathematical Theory}

\subsection{The Arthur--Selberg Trace Formula}

\begin{definition}[Automorphic Representation]
Let $G$ be a reductive group over a number field $F$. An \textbf{automorphic representation} of $G(\A_F)$ (where $\A_F$ is the adele ring) is an irreducible representation of $G(\A_F)$ occurring in the space $L^2(G(F)\backslash G(\A_F))$ of automorphic forms.
\end{definition}

\begin{theorem}[Arthur Trace Formula]
For a test function $f$ on $G(\A_F)$ with suitable properties, we have:
\[
\sum_{\pi} m_\pi \, \tr \pi(f) = \sum_{\{M\}} \frac{|\det(W_0^M)|}{|W_0^G|} \sum_{\gamma \in \mathfrak{C}(M)} a^M(\gamma) \int_{G_\gamma(\A_F) \backslash G(\A_F)} f(g^{-1}\gamma g) \, dg
\]
where the left side is the spectral expansion (sum over automorphic representations $\pi$) and the right side is the geometric expansion (sum over conjugacy classes $\gamma$).
\end{theorem}

\subsection{Arthur's Classification for Classical Groups}

\begin{definition}[Langlands Parameter]
A \textbf{global Langlands parameter} for $G$ is a homomorphism $\varphi: L_F \to {}^L G$ where $L_F$ is the Langlands group of $F$ and ${}^L G$ is the Langlands dual group.
\end{definition}

\begin{theorem}[Arthur's Endoscopic Classification, 2013]
Let $G$ be a classical group (symplectic, orthogonal, or unitary) over a number field $F$. There is a canonical bijection between:
\begin{enumerate}
    \item The set of global Langlands packets of square-integrable automorphic representations of $G$.
    \item The set of global Arthur parameters $\psi$ mapping into the dual group.
\end{enumerate}
This classification establishes that the discrete automorphic spectrum of $G$ is parametrized by certain self-dual automorphic representations of $\GL(N)$.
\end{theorem}

\subsection{The Arthur--Clozel Base Change}

\begin{theorem}[Arthur--Clozel, 1989]
Let $E/F$ be a cyclic extension of number fields of prime degree. Then there is a base change lifting from automorphic representations of $\GL(n, \A_F)$ to automorphic representations of $\GL(n, \A_E)$, realizing the Langlands functoriality for the inclusion $\GL(n,\C) \hookrightarrow \GL(n,\C)^d$ induced by the Galois group.
\end{theorem}

\subsection{The Local Trace Formula}

\begin{definition}[Local Trace Formula]
For a reductive group $G$ over a local field $F$, the \textbf{local trace formula} is an identity:
\[
\tr \pi(f) = \int_{G(F)} f(g) \Theta_\pi(g) \, dg
\]
where $\Theta_\pi$ is the character of the representation $\pi$.
\end{definition}

Arthur developed the local trace formula to relate characters of representations to orbital integrals, providing a powerful tool for the local Langlands correspondence.

\subsection{The Spectral and Geometric Sides}

The Arthur--Selberg trace formula admits a precise decomposition into two sides, each carrying distinct arithmetic and geometric information.

\begin{definition}[Spectral Side]
Let $G$ be a connected reductive group over a number field $F$, and let $f \in C_c^\infty(G(\A_F))$ be a smooth compactly supported test function. The \textbf{spectral side} of the trace formula is the expression
\[
\operatorname{Is}_{\mathrm{spec}}(f) = \sum_{\pi \in \Pi(G)} m(\pi) \cdot \operatorname{tr} \pi(f),
\]
where the sum runs over the set $\Pi(G)$ of automorphic representations of $G(\A_F)$, $m(\pi)$ denotes the multiplicity of $\pi$ in $L^2_0(G(F)\backslash G(\A_F))$ (the cuspidal subspace), and $\operatorname{tr}\pi(f)$ is the trace of $f$ acting on the representation space of $\pi$.
\end{definition}

The spectral side also includes contributions from the residual spectrum and the continuous spectrum, arising from the constant term of Eisenstein series. Explicitly, one has
\[
\operatorname{Is}_{\mathrm{spec}}(f) = \sum_{\pi \in \Pi_{\mathrm{cusp}}} m(\pi)\,\operatorname{tr}\pi(f) + \sum_{\mathfrak{a}} \int_{\mathfrak{a}^*_\C} \sum_{\pi \in \Pi_{\mathrm{res}}} m(\pi,\lambda)\,\operatorname{tr}\pi_\lambda(f)\,d\lambda,
\]
where $\Pi_{\mathrm{cusp}}$ denotes the cuspidal spectrum, $\Pi_{\mathrm{res}}$ the residual spectrum, and the second sum runs over relevant Levi subgroups $\mathfrak{a}$.

\begin{definition}[Geometric Side]
The \textbf{geometric side} of the trace formula is the expression
\[
\operatorname{Is}_{\mathrm{geom}}(f) = \sum_{\{M\}} \frac{|W_0^M|}{|W_0^G|} \sum_{\gamma \in \mathfrak{C}(M)} \int_{G_\gamma(\A_F)\backslash G(\A_F)} \Delta^G(\gamma, g)\, f(g^{-1}\gamma g)\,dg,
\]
where $\{M\}$ runs over the set of rational conjugacy classes of Levi subgroups of $G$, $W_0^M$ denotes the relative Weyl group of $M$ in $G$, $\mathfrak{C}(M)$ denotes the set of conjugacy classes of $M$, $G_\gamma$ is the centralizer of $\gamma$ in $G$, and $\Delta^G(\gamma, g)$ is a combinatorial weight depending on the structure of $G$ and the element $\gamma$.
\end{definition}

\begin{theorem}[Arthur--Selberg Trace Formula, Precise Form]
Let $G$ be a connected reductive group over a number field $F$, and let $f \in C_c^\infty(G(\A_F))$. Then
\[
\operatorname{Is}_{\mathrm{spec}}(f) = \operatorname{Is}_{\mathrm{geom}}(f),
\]
that is,
\[
\sum_{\pi \in \Pi(G)} m(\pi)\,\operatorname{tr}\pi(f) = \sum_{\{M\}} \frac{|W_0^M|}{|W_0^G|} \sum_{\gamma \in \mathfrak{C}(M)} \int_{G_\gamma(\A_F)\backslash G(\A_F)} \Delta^G(\gamma,g)\,f(g^{-1}\gamma g)\,dg.
\]
This identity holds provided that the test function $f$ is chosen to lie in an appropriate space of compactly supported functions satisfying certain regularity and support conditions (often denoted $C_c^\infty(G(\A_F))^{\sharp}$).
\end{theorem}

The key insight is that the spectral side encodes the automorphic spectrum (eigenvalues of operators on spaces of automorphic forms), while the geometric side encodes the conjugacy structure of $G$ over $F$. The trace formula thus provides a bridge between representation theory and arithmetic geometry.

\subsection{Stabilization of the Trace Formula}

The stabilization of the trace formula, due principally to Arthur, is the process of rewriting the geometric side in terms of stable orbital integrals rather than ordinary orbital integrals. This is essential for comparing trace formulas between different groups.

\begin{definition}[Stable Orbital Integral]
Let $\gamma \in G(F)$ be a semisimple element, and let $f \in C_c^\infty(G(\A_F))$. The \textbf{stable orbital integral} of $f$ over $\gamma$ is
\[
S\mathcal{O}_\gamma^G(f) = \int_{G_\gamma(\A_F)\backslash G(\A_F)} f(g^{-1}\gamma g)\,dg + \sum_{\gamma' \sim_s \gamma,\, \gamma' \neq \gamma} c(\gamma,\gamma')\,\int_{G_{\gamma'}(\A_F)\backslash G(\A_F)} f(g^{-1}\gamma' g)\,dg,
\]
where the sum runs over elements $\gamma'$ stably conjugate to but not rationally conjugate to $\gamma$, and $c(\gamma,\gamma')$ are explicit rational constants.
\end{definition}

\begin{remark}
The distinction between rational conjugacy and stable conjugacy is fundamental. Two elements $\gamma, \gamma' \in G(F)$ are \emph{stably conjugate} if $g^{-1}\gamma g = \gamma'$ for some $g \in G(\bar{F})$, whereas they are \emph{rationally conjugate} if $g \in G(F)$. Stable conjugacy is a coarser equivalence relation, and the stable orbital integral is invariant under stable conjugacy by construction.
\end{remark}

The geometric side of the trace formula decomposes into contributions from elliptic, hyperbolic, and unipotent orbital integrals. For the semisimple conjugacy classes:

\begin{itemize}
    \item \textbf{Elliptic orbital integrals:} These correspond to regular semisimple elements $\gamma$ for which the centralizer $G_\gamma$ is an anisotropic torus. These are the dominant terms and correspond to discrete series representations on the spectral side.
    \item \textbf{Hyperbolic orbital integrals:} These correspond to semisimple elements whose centralizer contains a split torus. They contribute to the continuous spectrum via the Eisenstein series.
    \item \textbf{Unipotent orbital integrals:} These correspond to unipotent elements $u \in G(F)$. The orbital integrals over unipotent conjugacy classes encode information about the residual spectrum and are computed using Harish-Chandra's unipotent integration formulas.
\end{itemize}

\begin{theorem}[Arthur's Stabilization]
There exists a stable test function $f^{\mathrm{st}} \in C_c^\infty(G(\A_F))$ such that the geometric side of the trace formula for $G$ can be rewritten as
\[
\operatorname{Is}_{\mathrm{geom}}(f) = \sum_{\{M\}} \frac{|W_0^M|}{|W_0^G|} \sum_{\gamma \in \mathfrak{S}(M)} S\mathcal{O}_\gamma^M(f^{\mathrm{st}}),
\]
where $\mathfrak{S}(M)$ denotes the set of stable conjugacy classes of $M$, and $S\mathcal{O}_\gamma^M$ denotes the stable orbital integral. Furthermore, the spectral side becomes a sum over stable packets of automorphic representations.
\end{theorem}

\begin{remark}
The stabilization requires control over the endoscopic contributions to the orbital integrals. Arthur showed that the weighted orbital integrals, which arise from his truncation procedure, can be decomposed into stable orbital integrals plus corrections from endoscopic groups. These corrections involve the so-called transfer factors $\Delta(\gamma, \delta)$, which are explicit functions relating conjugacy classes in $G$ to conjugacy classes in an endoscopic group $H$.
\end{remark}

\subsection{The Endoscopic Classification}

The endoscopic classification is one of the crowning achievements of Arthur's work and represents a major step toward the Langlands functoriality principle.

\begin{definition}[Arthur Parameter]
A \textbf{global Arthur parameter} for a reductive group $G$ over $F$ is a conjugacy class of homomorphisms
\[
\psi: L_F \times \SL_2(\C) \to {}^L G,
\]
where $L_F$ is the (conjectural) Langlands group of $F$ and ${}^L G$ is the Langlands dual group of $G$, such that:
\begin{enumerate}
    \item The restriction of $\psi$ to $L_F$ is a Langlands parameter (i.e., a homomorphism $L_F \to {}^L G$).
    \item The restriction of $\psi$ to $\SL_2(\C)$ factors through the principal $\SL_2$-embedding into ${}^L G$.
\end{enumerate}
\end{definition}

\begin{definition}[Tempered Representation]
An automorphic representation $\pi$ of $G(\A_F)$ is \textbf{tempered} if, for every parabolic subgroup $P = MN$ with Levi component $M$, the Jacquet module $\pi_N$ is a tempered representation of $M(\A_F)$. Equivalently, $\pi$ is tempered if and only if its Langlands parameters lie in the maximal compact subgroup of ${}^L G$.
\end{definition}

\begin{theorem}[Arthur's Endoscopic Classification, 2013]
Let $G$ be a classical group (symplectic $\mathrm{SO}(2n+1)$, orthogonal $\mathrm{SO}(2n)$ or $\mathrm{SO}(2n+1)$, or unitary $\mathrm{SU}(n)$) over a number field $F$. Then the discrete automorphic spectrum of $G$ is classified by Arthur parameters as follows:
\begin{enumerate}
    \item There is a canonical bijection between the set of global Arthur parameters $\psi$ for $G$ and the set of global Arthur packets $\Pi_\psi(G)$ of discrete automorphic representations.
    \item Each packet $\Pi_\psi(G)$ is parametrized by self-dual automorphic representations of $\GL(N, \A_F)$ for appropriate $N$.
    \item For tempered Arthur parameters (where $\psi|_{\SL_2(\C)}$ is trivial), the corresponding packets consist of tempered representations, and the classification reduces to the tempered endoscopic classification of Langlands and Rogawski.
\end{enumerate}
\end{theorem}

\begin{remark}
The proof of the endoscopic classification proceeds by comparison of the stable trace formula for $G$ with the trace formula for its endoscopic groups. The key ingredients are:
\begin{enumerate}
    \item \textbf{The Fundamental Lemma:} The identity relating orbital integrals on $G$ to orbital integrals on an endoscopic group $H$ for the test function arising from the unit element of the Hecke algebra. This was proved by Ng\^o and others.
    \item \textbf{Transfer of test functions:} A map from test functions on $G$ to test functions on $H$ that preserves the trace formula.
    \item \textbf{Local Langlands correspondence:} For the non-archimedean places, the local Langlands correspondence for classical groups (proved by Arthur, building on Harris--Taylor and Henniart for $\GL(n)$, and on Rogawski for unitary groups) is used to identify the local components of automorphic representations.
\end{enumerate}
\end{remark}

The endoscopic classification has profound consequences. It implies, among other things:
\begin{itemize}
    \item The existence of Langlands functorial transfer from classical groups to $\GL(n)$.
    \item The classification of supercuspidal representations of classical groups.
    \item The verification of the local Langlands conjectures for all classical groups over $p$-adic fields.
\end{itemize}

\subsection{The Trace Formula for \texorpdfstring{$\GL(n)$}{GL(n)}}

The case $G = \GL(n)$ occupies a special position in the Langlands program, as it is the group for which the Langlands correspondence is best understood and serves as the target for functorial transfers from other groups.

\begin{definition}[Test Function for $\GL(n)$]
For $G = \GL(n)$ over $F$, a \textbf{test function} is an element $f \in C_c^\infty(\GL(n, \A_F))$ of the form $f = \bigotimes_v f_v$, where $f_v \in C_c^\infty(\GL(n, F_v))$ for each place $v$ of $F$, and $f_v$ equals the characteristic function of $\GL(n, \mathcal{O}_v)$ for all but finitely many $v$.
\end{definition}

\begin{theorem}[Trace Formula for $\GL(n)$]
For $G = \GL(n)$ over a number field $F$, the trace formula takes the form:
\begin{align}
\sum_{\pi \in \Pi(\GL(n))} m(\pi)\,\operatorname{tr}\pi(f) &= \sum_{\gamma \in \mathfrak{C}(\GL(n))} \frac{1}{|\det(\Ad(\gamma) - 1)|_{\mathfrak{gl}(n)/\mathfrak{z}}|} \cdot \operatorname{vol}(\GL(n,F)\backslash \GL(n,\A_F))^{\mathrm{cusp}} \notag \\
&\quad \times \prod_v \int_{\GL(n,F_v)_\gamma \backslash \GL(n,F_v)} f_v(g^{-1}\gamma g)\,\frac{dg_v}{|\det(\Ad(\gamma)-1)|_{\mathfrak{gl}(n)/\mathfrak{z}}|^{1/2}}\,dg_v.
\end{align}
For $\GL(n)$, the geometric side simplifies because:
\begin{enumerate}
    \item Every semisimple element is elliptic (i.e., its centralizer is a torus), so the distinction between elliptic and hyperbolic orbital integrals becomes moot.
    \item The only Levi subgroups relevant to the cuspidal trace formula are the Borel subgroups.
    \item The transfer factors are explicit and well-understood.
\end{enumerate}
\end{theorem}

The trace formula for $\GL(n)$ plays a central role in Langlands functoriality in the following way. Let $H$ be another reductive group, and let ${}^L H \to \GL(n, \C)$ be a homomorphism of Langlands dual groups. Functoriality predicts a transfer map
\[
\mathfrak{F}: \Pi(H) \to \Pi(\GL(n)),
\]
and the trace formula comparison is the analytic vehicle for establishing such transfers. Specifically, one shows that for appropriate test functions $f$ on $H$ and $f^{\natural}$ on $\GL(n)$,
\[
\operatorname{Is}_{\mathrm{spec}}^H(f) = \operatorname{Is}_{\mathrm{spec}}^{\GL(n)}(f^{\natural}),
\]
which implies the transfer of automorphic representations from $H$ to $\GL(n)$.

\begin{remark}[Arthur's Function $f^\natural$]
For $\GL(n)$, the Arthur--Clozel base change and the endoscopic classification both rely on a specific construction of the transfer function $f^\natural$ on $\GL(n)$ associated to a test function $f$ on $G$. The function $f^\natural$ is defined by the identity
\[
S\mathcal{O}_\gamma^G(f) = S\mathcal{O}_{\mathfrak{F}(\gamma)}^{\GL(n)}(f^\natural),
\]
where $\mathfrak{F}(\gamma)$ denotes the image of the conjugacy class of $\gamma$ under the transfer map. The existence and uniqueness of such $f^\natural$ is one of the fundamental results in the theory.
\end{remark}

\begin{remark}
The trace formula for $\GL(n)$ is also the starting point for the \emph{relative trace formula}, which compares two different trace formulas on the same group. This technique, developed by Jacquet, Moeglin, and Waldspurger, has applications to the Gross--Prasad conjecture and the Ichino--Ikeda conjecture on central values of $L$-functions.
\end{remark}

\section{Impact on Science and Technology}

\subsection{In Mathematics}

\begin{enumerate}
    \item \textbf{Langlands Program:} Arthur's trace formula and endoscopic classification are central to the modern Langlands program, which connects number theory, representation theory, and harmonic analysis.
    \item \textbf{Number Theory:} Applications to the Birch and Swinnerton--Dyer conjecture, modularity of elliptic curves, and the Sato--Tate conjecture.
    \item \textbf{Representation Theory:} Arthur's conjectures guide the classification of representations of reductive groups over local and global fields.
    \item \textbf{Harmonic Analysis:} The trace formula is a powerful tool for studying the harmonic analysis of reductive groups.
\end{enumerate}

\subsection{In Other Fields}

\begin{enumerate}
    \item \textbf{Physics:} Automorphic forms and the trace formula appear in string theory (mirror symmetry, S-duality) and quantum chaos.
    \item \textbf{Cryptography:} The Langlands program has connections to isogeny-based cryptography (SIDH).
\end{enumerate}

\section{Frequently Asked Questions}

\subsection*{Q1: What is the Arthur--Selberg trace formula?}
A vast generalization of Selberg's trace formula that applies to any reductive group over a number field. It relates sums over automorphic representations (spectral side) to sums over conjugacy classes (geometric side) and is a central tool in the Langlands program.

\subsection*{Q2: What is Arthur's endoscopic classification?}
A complete description of the automorphic spectrum of classical groups in terms of automorphic representations of general linear groups. It establishes the Langlands functoriality principle for all classical groups.

\subsection*{Q3: What is the Langlands program?}
A far-reaching network of conjectures connecting number theory (Galois groups) to harmonic analysis (automorphic forms) through the Langlands dual group. It is one of the central themes of modern mathematics.

\subsection*{Q4: What is base change?}
Given an extension of number fields $E/F$ and an automorphic representation of $\GL(n)$ over $F$, base change produces an automorphic representation over $E$. The Arthur--Clozel theorem proves this for cyclic extensions.

\subsection*{Q5: What should an undergraduate study to understand Arthur's work?}
Number theory (algebraic and analytic), representation theory of Lie groups, harmonic analysis on groups, and algebraic geometry. Recommended: \textit{Automorphic Forms and Representations} (Bump), \textit{An Introduction to the Langlands Program} (Gelbart), and \textit{Trace Formulas} (Arthur--Elstrodt--Grunewald).

\section{Related Chapters}
See also Chapter~16 (Selberg) on trace formula and automorphic forms.

\section*{Exercises for the Reader}
\addcontentsline{toc}{section}{Exercises}

\begin{enumerate}
    \item [A] Derive Selberg's trace formula for a compact hyperbolic surface.
    \item [B] Show that the conjugacy classes in $\GL(n,\R)$ are classified by their Jordan normal form.
    \item [I] Prove that a non-trivial unipotent element cannot be a rational conjugacy class in $\GL(n,\Q)$.
    \item [I] Compute the Langlands dual group for $\SO(2n+1)$ and $\Sp(2n)$.
    \item [I] Show that the trace of an induced representation can be expressed in terms of orbital integrals.
    \item [B] State the Langlands functoriality conjecture for $\GL(2) \to \GL(3)$ via the symmetric square.
\end{enumerate}

\section*{Timeline}
\addcontentsline{toc}{section}{Timeline}

\begin{tabularx}{\linewidth}{|p{1.5cm}|>{\raggedright\arraybackslash}X|}
\hline
\textbf{Year} & \textbf{Event} \\ \hline
1944 & Born in Hamilton, Ontario, Canada \\ \hline
1966 & B.Sc.\ from the University of Toronto \\ \hline
1970 & Ph.D.\ from Yale (advisor: Robert Langlands) \\ \hline
1978 | Begins systematic development of the trace formula \\ \hline
1989 | Arthur--Clozel base change for $\GL(n)$ \\ \hline
1990 | Invariant trace formula \\ \hline
1991 | Professor at the University of Toronto \\ \hline
1999 | CRM-Fields-PIMS Prize \\ \hline
2005 | Steele Prize of the AMS (with Langlands) \\ \hline
2013 | Endoscopic classification published \\ \hline
2015 | Wolf Prize in Mathematics \\ \hline
2017 | Siberian Federal University Prize \\ \hline
\end{tabularx}

\section*{Glossary}
\addcontentsline{toc}{section}{Glossary}

\begin{description}
    \item[Automorphic form] A function on $G(F)\backslash G(\A_F)$ satisfying certain growth and invariance conditions.
    \item[Automorphic representation] An irreducible representation of $G(\A_F)$ in the space of automorphic forms.
    \item[Base change] A lift of automorphic representations from a group over $F$ to the same group over an extension $E$.
    \item[Endoscopy] A theory that relates the trace formula of a group to that of its endoscopic groups.
    \item[Langlands dual] A reductive group ${}^L G$ over $\C$ whose root datum is dual to that of $G$.
    \item[Langlands program] A web of conjectures connecting Galois theory, automorphic forms, and representation theory.
    \item[Orbital integral] An integral of a function over a conjugacy class.
    \item[Trace formula] An identity relating a sum over representations to a sum over conjugacy classes.
\end{description}

\section{Citations}\subsection*{Primary Sources}
\begin{enumerate}
    \item J.\ Arthur, \textit{A trace formula for reductive groups I}, Duke Math.\ J.\ 45 (1978), 911--952.
    \item J.\ Arthur, \textit{The invariant trace formula}, J.\ Amer.\ Math.\ Soc.\ 1 (1988), 323--383.
    \item J.\ Arthur, \textit{An introduction to the trace formula}, in \textit{Harmonic Analysis, the Trace Formula, and Shimura Varieties}, AMS, 2005.
    \item J.\ Arthur, \textit{The endoscopic classification of representations}, AMS Colloquium Publ., 2013.
    \item J.\ Arthur and L.\ Clozel, \textit{Simple Algebras, Base Change, and the Advanced Theory of the Trace Formula}, Ann.\ Math.\ Studies 120, Princeton Univ.\ Press, 1989.
\end{enumerate}

\subsection*{Expository Works}
\begin{enumerate}
    \item J.\ Arthur, \textit{The principle of functoriality}, Bull.\ AMS 40 (2003), 39--53.
    \item S.\ Gelbart, \textit{An elementary introduction to the Langlands program}, Bull.\ AMS 10 (1984), 177--219.
    \item R.\ Langlands, \textit{Problems in the theory of automorphic forms}, Lecture Notes in Math.\ 170, Springer, 1970.
\end{enumerate}

